\subsubsection{User Interface}
Zur Planung der Benutzeroberfläche wurden im Vorfeld Wireframes und Mockups mit dem Design-Tool Figma erstellt. 
Diese dienen dazu, die Struktur der einzelnen Seiten frühzeitig zu visualisieren und zu überprüfen. 
Interaktive Funktionalitäten oder Navigationsflüsse wurden im Rahmen dieses Projekts nicht prototypisch umgesetzt. 
Die Entwürfe umfassen unter anderem:
\begin{itemize}
	\item die Karriereseite
	\item die Stellenübersicht
	\item die detailierte Stellenanzeige
	\item ein Bewerbungsformular für Kandidat:innen
	\item ein Dashboard fuer Recruiter:innen  
\end{itemize}
Die wichtigsten Benutzeroberflächen sind in den folgenden Abbildungen dargestellt.
\begin{figure}[H]
    \centering
    
\includegraphics[width=0.7\textwidth]{LandingPage.png}
    \caption{Mockup der Landing Page.}
    \label{fig:landingPage}
\end{figure}
Die Seite dient als Einstiegspunkt für Bewerber:innen, stellt allgemeine Unternehmensinformationen dar und bietet eine Weiterleitung zu den offenen Stellenanzeigen.
\begin{figure}[H]
    \centering
    \includegraphics[width=0.7\textwidth]{StellenOverview.png}
    \caption{Mockup der Stellenübersicht.}
    \label{fig:stellenOverview}
\end{figure}
Bewerber:innen können hier offene Positionen einsehen und über eine Detailansicht weitere Informationen zur jeweiligen Stelle abrufen.
\begin{figure}[H]
    \centering
    \includegraphics[width=0.7\textwidth]{Stellenanzeige.png}
    \caption{Mockup der Stellendetailseite.}
    \label{fig:stelle}
\end{figure}
Diese Seite enthält detaillierte Informationen zur ausgeschriebenen Position sowie eine Möglichkeit, den Bewerbungsprozess zu starten.
\begin{figure}[H]
    \centering
    \includegraphics[width=0.7\textwidth]{Bewerbungsformular.png}
    \caption{Wireframe des Bewerbungsformulars.}
    \label{fig:bewerbungFormular}
\end{figure}
Über dieses Formular können Bewerber:innen ihre persönlichen Daten eingeben und Bewerbungsunterlagen wie Lebenslauf oder Anschreiben hochladen.
\begin{figure}[H]
    \centering
    \includegraphics[width=0.7\textwidth]{RecruiterDashboard.png}
    \caption{Wireframe des Recruiter-Dashboards.}
    \label{fig:recruiterDashboard}
\end{figure}
Recruiter:innen können hier Stellenanzeigen verwalten sowie eingegangene Bewerbungen einsehen und weiterfuehrend bearbeiten.
\newline
Die weiteren erstellten WireFrames bzw. Mockups koennen unter folgenden 
\href{https://www.figma.com/design/0LDsbppkYlQgEbpulk5ud9/ATS-WebApp?node-id=20-46&t=cHVE2zw3fKD8hl14}{\textit{Link}} eingesehen werden.
